\providecommand{\documentfontsize}{12pt}
\providecommand{\documentpaper}{letterpaper}
\edef\documentclassoptions{\documentfontsize,\documentpaper}
\expandafter\documentclass\expandafter[\documentclassoptions]{article}
\usepackage{document}
\addbibresource{references.bib}

% Only the example uses PGFPlots; the style accepts ordinary image files too.
\usepackage{pgfplots}
\pgfplotsset{compat=1.18}

\title{A Minimalist Research Document\thanks{%
  Email: \href{mailto:author@example.org}{author@example.org}.
  For helpful comments we thank xxx. All potential errors are our own.}}
\author{Yuntian Zhao\\
  {\small\color{GitHubMuted}Research document template}}
\date{\today}
\hypersetup{
  pdftitle={A Minimalist Research Document},
  pdfauthor={Yuntian Zhao},
  pdfsubject={Research article template with Libertinus Serif}
}

\begin{document}
\maketitle

\begin{abstract}
This template adapts the typography and restrained visual language of the
companion Beamer template to a research article. Libertinus Serif provides
body text and headings, with matching mathematical typography and a separate
monospaced code font. A quadratic optimization problem illustrates equations,
assumptions, propositions, proofs, figures, tables, and algorithms. The examples
are self-contained and serve only to demonstrate document structure.
\keywords{research writing; mathematical typography; reproducible documents}
\end{abstract}

\section{Introduction}\label{sec:introduction}

A research document should make its argument easy to follow without making
its visual design the subject of attention. This template keeps the Beamer
version's dark foreground, blue links, muted notes, and light code backgrounds.
Plain bold titles and section headings accompany a single-column article,
one-inch margins, page numbers, and one-and-a-half-spaced prose.

Content belongs in the \texttt{.tex} file; reusable formatting belongs in
\texttt{document.sty}. Start a new paper from \texttt{document-minimal.tex}.
The full example shows common research-writing patterns without requiring
external figures or code files. Standard \LaTeX{} commands remain available,
including \verb|\section|, \verb|\label|, \verb|\ref|, and \verb|\appendix|.

Author--date citations use the same Chicago style as the slide template.
A textual citation looks like \textcite{melitz_impact_2003}; a parenthetical
citation looks like \parencite{melitz_impact_2003}. The bibliography database
is declared in the document rather than in the style, so each paper can use
its own reference collection. The cited article is an example of the citation
format, not a source for the optimization example below.

\section{A running example}\label{sec:model}

\subsection{Setup and notation}

Let $\bm{\theta}\in\mathbb{R}^d$ denote the parameter vector and let
$\bm{\theta}^{\star}$ denote its optimal value. Consider the quadratic objective
\begin{equation}\label{eq:objective}
  f(\bm{\theta})
  = \frac{1}{2}(\bm{\theta}-\bm{\theta}^{\star})^{\mathsf{T}}
    Q(\bm{\theta}-\bm{\theta}^{\star}).
\end{equation}
The expression separates parameters from curvature and gives a compact example
of vectors, superscripts, and equation numbering. References such as
Equation~\eqref{eq:objective} remain linked in the PDF.

\begin{assumption}[Positive curvature]\label{ass:curvature}
The matrix $Q\in\mathbb{R}^{d\times d}$ is symmetric and positive definite,
with eigenvalues in $[\mu,L]$, where $0<\mu\leq L$.
\end{assumption}

\begin{definition}[Gradient update]
For a fixed step size $\eta>0$, the sequence of iterates is defined by
\begin{equation}\label{eq:update}
  \bm{\theta}_{k+1}
  = \bm{\theta}_k-\eta\nabla f(\bm{\theta}_k),
  \qquad k=0,1,2,\ldots.
\end{equation}
\end{definition}

\subsection{A convergence result}

\begin{proposition}[Contraction]\label{prop:contraction}
Under Assumption~\ref{ass:curvature}, if $0<\eta<2/L$, then
\begin{equation}\label{eq:contraction}
  \lVert\bm{\theta}_{k}-\bm{\theta}^{\star}\rVert_2
  \leq q^k\lVert\bm{\theta}_{0}-\bm{\theta}^{\star}\rVert_2,
  \qquad q=\max_{\lambda\in\{\mu,L\}}|1-\eta\lambda|<1.
\end{equation}
\end{proposition}

\begin{proof}
The gradient is $Q(\bm{\theta}-\bm{\theta}^{\star})$. Subtracting
$\bm{\theta}^{\star}$ from both sides of Equation~\eqref{eq:update} gives
\[
  \bm{\theta}_{k+1}-\bm{\theta}^{\star}
  =(I-\eta Q)(\bm{\theta}_{k}-\bm{\theta}^{\star}).
\]
By symmetry of $Q$, the spectral norm of $I-\eta Q$ is bounded above by
$q$. The step-size restriction makes both endpoint magnitudes strictly less
than one. Iterating the norm inequality proves the claim.
\end{proof}

\begin{remark}
Theorem-like environments share a counter within each section. Assumptions
and definitions use upright text; propositions and theorems use italic text.
Proofs use the standard end-of-proof square.
\end{remark}

\section{Illustrations and computation}\label{sec:illustrations}

\subsection{Figures and tables}

Figure~\ref{fig:objective} shows a scalar instance of
Equation~\eqref{eq:objective}. Captions describe the object; notes supply
assumptions, sources, and qualifications in muted text. For an external figure,
replace the plot with:
\begin{lstlisting}[language={[LaTeX]TeX},numbers=none]
\includegraphics[width=\linewidth]{your-figure}
\end{lstlisting}

\begin{figure}[tbp]
  \centering
  \begin{tikzpicture}
    \begin{axis}[
      width=0.78\linewidth,height=0.37\linewidth,
      xmin=-2,xmax=2,ymin=0,ymax=2.2,
      xlabel={$\theta$},ylabel={$f(\theta)$},
      axis lines=left,
      axis line style={GitHubMuted},
      tick label style={font=\small},
      label style={font=\small},
      xtick={-2,-1,0,1,2},ytick={0,1,2},
      samples=81,domain=-2:2]
      \addplot[GitHubBlue,thick] {0.5*x^2};
      \addplot[only marks,mark=*,GitHubRed] coordinates {(0,0)};
    \end{axis}
  \end{tikzpicture}
  \caption{A quadratic objective and its minimizer}\label{fig:objective}
  \figurenote{\textit{Notes:} The scalar example sets $Q=1$ and
    $\theta^{\star}=0$. The dot marks the unique minimizer.}
\end{figure}

Table~\ref{tab:convergence} illustrates the geometric bound in
Proposition~\ref{prop:contraction} for two possible contraction factors.
Horizontal rules separate the header and body; vertical rules are omitted.
The numbers are analytical illustrations rather than empirical estimates.

\begin{table}[tbp]
  \centering
  \caption{Illustrative normalized error bounds}\label{tab:convergence}
  \begin{tabular}{@{}rrr@{}}
    \toprule
    Iteration $k$ & $q=0.5$ & $q=0.8$ \\
    \midrule
     0 & $1.0000$ & $1.0000$ \\
     5 & $0.0313$ & $0.3277$ \\
    10 & $0.0010$ & $0.1074$ \\
    20 & $9.54\times10^{-7}$ & $0.0115$ \\
    \bottomrule
  \end{tabular}
  \figurenote{\textit{Notes:} Each entry equals $q^k$, rounded for display.
    The initial error is normalized to one.}
\end{table}

\subsection{Algorithms and code}

Algorithm~\ref{alg:gradient} states the update independently of any programming
language. A fixed iteration count keeps the example's stopping rule explicit.
In an application, report the step size, initialization, and convergence
criterion alongside the results.

\begin{algorithm}[tbp]
  \caption{Gradient descent with a fixed iteration count}\label{alg:gradient}
  \begin{algorithmic}[1]
    \Require Objective $f$, initial value $\bm{\theta}_0$, step size $\eta$,
      iteration count $K$
    \For{$k=0,\ldots,K-1$}
      \State $\bm{\theta}_{k+1}\gets\bm{\theta}_k-\eta\nabla f(\bm{\theta}_k)$
    \EndFor
    \State \Return $\bm{\theta}_K$
  \end{algorithmic}
\end{algorithm}

Listing~\ref{lst:gradient} implements the scalar case using only standard
Python. Cascadia Code keeps operators and punctuation distinct, with code
ligatures disabled. Source files can instead be included with:
\begin{lstlisting}[language={[LaTeX]TeX},numbers=none]
\lstinputlisting[language=Python]{your-code.py}
\end{lstlisting}

\begin{lstlisting}[language=Python,caption={Gradient descent for the scalar objective},label={lst:gradient}]
def gradient_descent(theta: float, step: float,
                     iterations: int) -> float:
    """Minimize f(theta) = theta**2 / 2 for 0 < step < 2."""
    for _ in range(iterations):
        theta -= step * theta
    return theta

estimate = gradient_descent(1.0, 0.5, 20)
assert abs(estimate) <= 1e-6
\end{lstlisting}

\subsection{Lists and emphasis}

Use lists to isolate a small number of distinct conditions or conclusions:
\begin{itemize}
  \item Identify the estimand or theoretical object before presenting results.
  \item State the assumptions needed for the argument.
    \begin{itemize}
      \item Distinguish maintained assumptions from testable implications.
    \end{itemize}
  \item Keep qualifications close to the claim they qualify.
\end{itemize}
For ordered steps, use an enumeration:
\begin{enumerate}
  \item Specify the model or empirical design.
  \item Present the result and its uncertainty.
  \item Discuss robustness and external validity.
\end{enumerate}
Use \emph{italics} for emphasis and \textbf{bold text} sparingly. Footnotes
remain available for details that would interrupt the argument.\footnote{The
body is one-and-a-half spaced; footnotes, captions, code, and bibliography
entries use compact spacing.}

\section{Conclusion}

The template preserves the presentation's visual identity while using the
conventions of a paper: serif body text, numbered sections, floating figures
and tables, cross-references, an author--date bibliography, and lettered
appendices. Keep the visual defaults in the style file and adapt the content
to the research question.

\printbibliography[heading=bibintoc,title={References}]

\clearpage
\appendix
\section{Mathematical notation}\label{app:notation}

The appendix uses letters for section numbers. Equation and theorem numbering
follow the section automatically, so appendix material has distinct labels:
\begin{equation}\label{eq:appendix}
  \mathbb{E}[X]=\int_{\Omega}X(\omega)\,\mathrm{d}\mathbb{P}(\omega).
\end{equation}
The following alphabets retain the distinctions used by the Beamer template:
\begin{align*}
  \text{Calligraphic:}\quad & \mathcal{ABC} && \mathbcal{ABC} \\
  \text{Script:}\quad       & \mathscr{ABC} && \mathbscr{ABC} \\
  \text{Fraktur:}\quad      & \mathfrak{ABC} && \mathbfrak{ABC} \\
  \text{Blackboard:}\quad   & \mathbb{R}\quad\mathbb{E}\quad\mathbb{P}.
\end{align*}

\begin{lemma}[Scalar iteration]
If $d=1$, $Q=1$, and $\theta^{\star}=0$, then
$\theta_k=(1-\eta)^k\theta_0$.
\end{lemma}
\begin{proof}
Substitute the scalar gradient into Equation~\eqref{eq:update} and iterate.
\end{proof}

\section{Long-table pattern}

Use \texttt{longtable} directly, rather than nesting it inside a floating table.
Its header and footer definitions repeat automatically when rows cross a page
boundary. The short glossary below supplies the pattern for longer data tables.

\begin{longtable}{@{}p{0.18\linewidth}p{0.73\linewidth}@{}}
  \caption{Notation for the running example}\label{tab:notation}\\
  \toprule
  Symbol & Meaning \\
  \midrule
  \endfirsthead
  \multicolumn{2}{l}{\small\tablename~\thetable\ (continued)}\\
  \toprule
  Symbol & Meaning \\
  \midrule
  \endhead
  \midrule
  \multicolumn{2}{r}{\small Continued on the next page}\\
  \endfoot
  \bottomrule
  \endlastfoot
  $d$ & Dimension of the parameter vector. \\
  $\bm{\theta}$ & Parameter vector. \\
  $\bm{\theta}^{\star}$ & Unique minimizer of the objective. \\
  $Q$ & Symmetric positive-definite curvature matrix. \\
  $\mu$ & Positive lower bound on the eigenvalues of $Q$. \\
  $L$ & Upper bound on the eigenvalues of $Q$. \\
  $\eta$ & Step size in the gradient update. \\
  $q$ & Upper bound on the per-iteration contraction factor. \\
  $k$ & Iteration index, starting at zero. \\
  $K$ & Number of gradient updates performed. \\
\end{longtable}

\end{document}
